\documentclass[../main.tex]{subfiles}
\begin{document}
\onlyinsubfile{
\setcounter{chapter}{0}
}
\notinsubfile{}

\section{De module in een oogopslag}\label{chap:Hoog}
%\marginpar{\nogdoen{module-cie moet dit hier of in docentenhandleiding? Tekst is gekopi\"eerd uit par 1.1}}
In deze module maak je kennis met quantumcomputing. De module bestaat uit twee delen. In het eerste deel bouwen we ons begrippenapparaat op, in het tweede deel passen we onze kennis toe in praktische opdrachten. Het eerste deel bestaat uit drie hoofdstukken. Alle hoofdstukken hebben een aantal werkbladen. In hoofdstuk~\ref{chap:H1} beginnen we met een korte inleiding op het onderwerp en maken we kennis met enkele belangrijke begrippen en fenomenen door zelf te experimenteren met licht. Om de vergelijking te maken met een quantumcomputer begint hoofdstuk~\ref{chap:H2} met een korte uitleg over hoe een klassieke computer werkt. We onderbouwen ons wiskundig begrippenapparaat met een model. Het hoofdstuk eindigt met een eisenlijst waaraan een quantumcomputer moet voldoen. In ~\ref{chap:H3} behandelen we de bouwblokken (poorten) van quantumcomputers. We zien dat die echt anders zijn. We leren hiermee werken en het hoofdstuk sluit af met een protocol waar alle eigenschappen van een quantumcomputer in voorkomen: teleportatie. Na dit hoofdstuk volgt een toets. 

Het tweede deel (hoofdstuk~\ref{chap:H4}) is een verzameling praktische keuze-opdrachten. Daarin leer je pas goed hoe de quantumcomputer is toe te passen. Je onderzoekt in samenwerkingsverband een aspect van quantumcomputing; een algoritme, een verdieping in wiskunde D, een technisch of een maatschappelijk aspect. Er is ook de mogelijkheid om dit deel in te vullen met een activiteit aangeboden door een van de universiteiten. 
\iffalse
het tweede deel bestaat uit keuzeopdrachten waarvoor je in de lessen aanwijzingen krijgt. In hoofdstuk~\ref{chap:H4} werken we met quantumalgoritmen. Met spelletjes maken we nog eens duidelijk dat een quantumcomputer echt andere dingen kan dan een klassieke computer. Hoofdstuk~\ref{chap:H5} behandelt hoe je een quantumcomputer bouwt. Hoofdstuk~\ref{chap:H6} schetst aan de hand van enkele concrete voorbeelden welke problemen quantumcomputers een rol kunnen spelen in de maatschappij en we speculeren we voorzichtig over de toekomst. Na hoofdstuk zes volgen nog verdiepingsopdrachten voor wiskunde~D. In werkbladen staan enkele praktische opdrachten uitgewerkt. Een uitgebreide literatuurlijst met voornamelijk open access bronnen geeft je toegang tot oorspronkelijke literatuur. 
\fi
%\bigskip
\marginpar{\vspace{-4cm}\includegraphics[width=0.95\marginparwidth]{./img/smartphone verplicht.png}
Tip: Voor de papieren versie zijn de links als QR-codes opgenomen. Gebruik je telefoon!}

Deze reader is het best te lezen in elektronisch (pdf) formaat. Voor de papieren versie zijn voor de links QR-codes in de marge opgenomen. We kunnen niet al ons materiaal in deze reader kwijt. In de docentenhandleiding en op de~\hrefqr{www.quantumrules.nl}{website} van het project staat meer ondersteunend materiaal.

\bigskip
Veel plezier,

De auteurs
\end{document}